% =====================================================================
%  Harald Høffding, Montaigne's Significance in the History of Ethics
%  Translated from the Danish: „Montaignes Betydning i Etikens Historie“, Det nittende Aarhundrede (April--September 1876), June 1876, pp. 210--230.
%
%  TRANSLATION -- GENERATED; DO NOT EDIT BY HAND.  Built by ebook-method/hoeffding-1875-76/tr_parts.py build montaigne-1876 --write
%  from the parts in ebook-method/hoeffding-1875-76/tr/ (GLOSSARY.md, TRANSLATOR-BRIEF.md).  Made from the collated
%  transcription texts/hoeffding/montaigne-1876/transcription.tex (2026-10-15); [N] in the margin = the page of the original.
%  Terms: GLOSSARY.md (choices marked ◆ there for Hans to review).
% =====================================================================
\documentclass[12pt,a4paper,openany]{book}
\usepackage[utf8]{inputenc}
\usepackage[T1]{fontenc}
\usepackage{libertinus}
\usepackage{libertinust1math}
\usepackage[german,french,english]{babel}
\usepackage[protrusion=true,expansion=true]{microtype}
\usepackage{setspace}
\usepackage[margin=1.3in]{geometry}
\usepackage{fancyhdr}
\usepackage{hyperref}
\hypersetup{pdftitle={Montaigne's Significance in the History of Ethics}, pdfauthor={Harald Høffding}, hidelinks}
\pagestyle{fancy}
\fancyhf{}
\fancyhead[LE,RO]{\thepage}
\fancyhead[C]{\textit{H. Høffding: Montaigne's Significance in the History of Ethics (1876)}}
\fancypagestyle{plain}{\fancyhf{}\fancyhead[LE,RO]{\thepage}}
\setstretch{1.3}
\usepackage{marginnote}
\usepackage{graphicx}
\DeclareUnicodeCharacter{0254}{\reflectbox{c}}   % ɔ (id est)
\renewcommand{\thefootnote}{\ifcase\value{footnote}\or *\or **\or ***\else\arabic{footnote}\fi)}
\newcommand{\opage}[1]{\marginnote{\footnotesize\textit{[#1]}}}
\newcommand{\signature}[1]{\par\smallskip\hfill #1\hspace*{2em}\par}
\pdfstringdefDisableCommands{\renewcommand{\footnote}[2][]{}\renewcommand{\opage}[1]{}}

\begin{document}

\begin{titlepage}
\begin{center}
{\LARGE\bfseries The Nineteenth Century.}\\[0.6em]{\small A Monthly for Literature and Criticism}\\[0.4em]{\small edited by}\\[0.2em]{\large Georg Brandes and Edvard Brandes.}\\[0.8em]{\small 1876. April--September.}\\\vfill{\small Copenhagen. Published by Gyldendal (F. Hegel). 1876.}\\[2em]
{\small Translated from the Danish}
\end{center}
\end{titlepage}
\clearpage

% --- p. 210 ---
\setcounter{footnote}{0}%
\opage{210}\begin{center}{\Large\sffamily\bfseries Montaigne's Significance in the History of Ethics.}\\[0.5em]\rule{4em}{0.4pt}\end{center}\medskip

A parallel may be drawn between the rebirth of science in the 15th and 16th centuries and the turning point that
occurred in the history of Greek thought with Socrates. The Greek philosophers before him
had immersed themselves in the contemplation of the world, of external
nature, which the Greeks called the heavens. Socrates now insisted
that if man is to have a knowledge
that can help him in life and become the soul of his
action, he must first and foremost know himself.
The human, the nature, tasks and duties of man,
were for Socrates the essential object of thought, and
he thereby became the founder of ethics, the doctrine
of the conditions and aims of human action. In
relation to the earlier thinkers he stood, then, as the one
who brought philosophy down from the heavens to the houses
and dwellings of men. This image can be used also of the thinkers of the Renaissance. But the heavens then signify not,
as with the Greeks, external nature, the universe, but the
supernatural world, in which the medieval Church
had taught men to see the true goal of life, in
relation to which earthly life itself was a nothing.
Worldly life had no independent significance, no
essential tasks to solve; it was a short time of probation,
which was concluded by an absolute decision. The rebirth of art and
science was a sign that men were beginning once more
% --- p. 211 ---
\setcounter{footnote}{0}%
\opage{211}to turn back from the heavens to the human world.
They drew strength above all from the contemplation
and appropriation of ancient culture, in which they found the
purely human life mirrored, the ``natural man,''
whom the Church condemned. Nature as a whole became the object
of enthusiastic contemplation, and the flowering natural science
soon gave this contemplation a firm and unshakable
foundation. Men no longer felt strangers in
the universe; there arose a confidence and a trust in
life that made a striving toward the solution of purely
human tasks possible. But this trust showed itself,
as a consequence of the whole character of the age, in an immediate,
tumultuous way, as burning enthusiasm, not
with clear self-possession. Man's
essence is still seen as interwoven in the connection of nature, and is not grasped
in its whole depth and inwardness. But Nature
stands as the infinite, all-embracing, in which all finite
oppositions are lost, above all the opposition of antiquity
between form and matter, the opposition of the Middle Ages between
mind and nature.

As regards ethics in particular, it now brought forward,
on the one hand, what the Church called the natural
man, with his drives and moods, his thoughts
and powers, as the real foundation for enjoyment and
action (Montaigne, Macchiavelli) or in order to show that
he has in himself the principle of moral action
and need not wait for the reward in a beyond (Pomponatius);
on the other hand, liberation is found in
devotion to and enthusiastic self-sacrifice for the infinite
content of ideas that now dawns upon consciousness (Giordano
Bruno). Naturally these two sides cannot be separated;
they stand in inner connection with one another; this follows
already from the preceding remarks, and it will
become evident above all in the following portrayal of one of the
thinkers of the Renaissance mentioned.

% --- p. 212 ---
\setcounter{footnote}{0}%
\opage{212}Montaigne is a typical representative of the skepticism
that, in times full of strong oppositions of mind,
often arises in clear, unprejudiced minds, which do not
feel the need of a positive view, if only for the
reason that one would then have to take the field for such a view
and thus become oneself one of the many contending parties.
This Montaigne would avoid at any price. He lived
in France (1533---1592) in the time of the wars of religion
and was here a witness to how each of the contending
parties not only believed that it had right on its
side, but was sure that it possessed the one saving
truth, so sure that it would bring its opponents over to its side with fire and sword.
Montaigne's
sound and independent nature --- he has been called
the most original mind of the 16th century --- developed
by a sensible and natural education, by travel and
by extensive reading, especially of classical authors,
could not but feel itself prompted to inquire into
the essence of human knowledge and conviction,
which could allow of such diversity and
strife. Yet one must not here expect to find methodical
investigations. Thinking is for him rather a
kind of play or pastime; he lets his thoughts go
as they will, and in turn gives his readers them just as
unbound and disordered as he himself got them\footnote[1]{``Je n'ay pas aultre sergeant de bande à renger mes pieces que la fortune: à mesme que mes resveries se presentent, je les entasse''. Essais de Michel de Montaigne liv. II chap. 10.}.
Hence he does not promise his readers learning either; the
only thing they will come to know from his book is himself:
C'est icy un livre de bonne foy, lecteur. Il t'advertit dez
l'entré, que je ne m'y suis proposé aulcune fin, que domestique
et privée''. (L'aucteur au lecteur),

But even this is significant. To give oneself up thus
to one's own thoughts and to rock oneself in one's
% --- p. 213 ---
\setcounter{footnote}{0}%
\opage{213}own fancies is a modern trait, a testimony that
the individual has already freed himself from the traditional
presuppositions and, without troubling himself about any authority,
confidently dares to surrender himself to the currents of his own
inner life and let them lead him on voyages of discovery to
new, unknown regions. This gives us reason already
to ask about Montaigne's relation to
the teaching of the Church.

In his youth Montaigne had, at his father's
request, translated a work by Raimundus of Sabunde
(who lived in the 15th century), a natural theology
that sought in nature and reason proofs of the
doctrines of the Church. Later he heard that from various
quarters objections had been raised against this argument,
and in the longest chapter among his ``Essays''
(liv. II chap. 12) he then sought to meet these objections.
But he does this in a manner that, despite
the sustained irony, betrays his real view.

On the one hand it had been held that the doctrines of faith
were too holy and exalted to be supported by
human grounds; on the other hand it had
declared the grounds employed to be weak and useless.

Toward the first opponents he shows himself seemingly
very benevolent. Very well, he says. Faith ought
to be enough in itself. But --- if our faith were really
genuine, inspired from heaven, then our life would surely look
different; all relations would be transformed in
accordance with the holy demands, and Christians
would not so often be put to shame in their
conduct by Turks and heathens. Nor would
the wars of religion, with all their ignoble passions, have
been able to split us; human affect and interest
would not have held the floor in these holy matters.
All this is proof enough that we always take religion
in our own way, with our human hands, and
that supernatural inspiration gives way to tradition,
% --- p. 214 ---
\setcounter{footnote}{0}%
\opage{214}custom and habit. We are Christians for similar reasons as
we are French or German.

Already here Montaigne has forgotten Raimundus, whom
he was to defend. But he loses sight of him entirely
as soon as he takes up the second
objection, that the proofs fall short. These opponents
he will crush completely by showing them the
impotence and nothingness of human knowledge. But
instead of the conclusion one would expect from this, that
man must blindly surrender to the divine
authority, it shines clearly through that Montaigne himself
accepts the skeptical results, and that for him they
form a limit not only against philosophical but
also against theological dogmatism. He raises
the Divine so high above man's sight and
power\footnote[1]{C'est aux chrestiens une occasion de croire, que de rencontrer chose incroyable. Liv. II. chap 10.}, not in order to worship it the more inwardly, but in order
to turn away from it with irony.

Montaigne sets the teachings of the various philosophical schools
against one another; he shows the unreliability of the senses;
he gathers traits from the life of animals to establish
that man possesses no essential advantage
over them; --- in short, he has the whole customary arsenal of
skepticism at his disposal, and he
uses it with virtuosity. But what ought especially to
interest us is to see how he knows how to concentrate
the skeptical arguments in a way that
betrays a genuine philosophical thinker behind the playful ironist.
Our reason, he says, cannot show us the truth;
for in order that we might be sure that it really
was the truth it showed us, we should need another
reason, which could vouch for the truth of the first, and
for this second one again a third, and so on to
infinity. All our knowledge ultimately stems from
% --- p. 215 ---
\setcounter{footnote}{0}%
\opage{215}the senses\footnote[1]{Toute cognoissance s'achemine en nous par les sens; se sont nos maistres (ibid).}, but neither can they teach us what is true.
For they apprehend the thing only as it appears to them
by virtue of their constitution\footnote[2]{Les sens ne comprennent pas le subject estrangier, ains seulement leurs propres passions (ibid).}. The thing itself eludes
us constantly, and we have no standard by which to
judge the relation between its real essence and our
sensation of it. In a similar way Montaigne shows
that, since the imagination cannot go beyond
our own essence, we necessarily think of the Deity
in human forms. Although men have thus
themselves fashioned their gods, they nevertheless become afraid
of them, just as children become afraid of the face of their
playmate which they have themselves painted or blackened. ---
There lies at the base of this an intimation of the thought which Kant
and Feuerbach in more recent times have carried through so energetically
in the domain of the theory of knowledge and the philosophy of religion,
that the seen is always conditioned by the sight and is not known
except in its relation to it.

Montaigne draws from this both theoretical and, above all,
practical consequences. Since opinion stands against opinion
among the learned, --- since there is nothing so absurd
that some people or other has not taken it up into its religion,
--- since the good and the right is different among the different
nations, so that what is sin on this side of
the mountains is virtue on the other, --- there is
nothing for it but to admit one's ignorance. It is
a proof of madness and of the very highest degree of
ignorance to regard oneself as certain of anything. All our
inquiry can only serve to develop and confirm the
ignorance that is innate in us. Indeed, one cannot even
be sure of one's doubt and ignorance; one should
therefore not say: I doubt! --- but: what do I know? ---
% --- p. 216 ---
\setcounter{footnote}{0}%
\opage{216}From this it follows that one may well excuse those who err in
their religious ideas; for since these ideas
concern something that lies beyond the range of human
grounds, one who is not enlightened by divine
grace may easily go astray; --- but that, on the
other hand, it is indefensible to set oneself against the existing
state of things because one believes oneself to have found
something that is better. That bespeaks great self-love
and partiality. I loathe the new, says Montaigne,
in whatever shape it presents itself. Not
because the existing laws are in themselves always reasonable.
The validity of laws rests not on their being just,
but precisely on their being laws: c'est le fondement mystique
de leur auctorité! The origin of laws often derives
from folly and self-interest, but custom
gives them firmness and makes them seem just. Therefore
it is dangerous to follow them back to their origin.
It is with them as with a mighty river,
which at its source is only an insignificant spring. Thus
the power of laws grows from a chance beginning in
the course of time.

Custom is the mistress of the world. This holds
not only of the juridical laws but also of the moral ones.
Conscience feels itself bound to what it is accustomed
to respect, and cannot free itself from it without repentance and rancor.
Cannibals, who regard it as a
duty of piety to eat their aged parents, would be shocked
to see the departed buried in the earth as
food for the worms. The wise man must indeed inwardly free
his soul, as much as possible withdraw from the crowd
in order to judge things without prejudice (retirer son
ame de la presse et la tenir en liberté et puissance de
juger librement des choses. Liv. I chap. 22); --- but outwardly
he must follow the existing laws and customs. In
contrast to the many conflicting opinions he surrenders
himself immediately and naively to Nature, which never
% --- p. 217 ---
\setcounter{footnote}{0}%
\opage{217}errs, so long as one does not falsify it by wanting to
interpret its wisdom; but against this precisely ignorance
and indifference secure him: ``Oh! que c'est un doulx
et mol chevet, et sain, que l'ignorance et l'incuriosité, à
reposer une teste bien faicte!'' (liv. III chap. 13). From
external occupations he withdraws into his own self
and follows it into its innermost secret passages, its darkest
folds. But this way to peace cannot then naturally
be common to all; each must seek it for himself, according to
his nature.

This thoroughgoing skepticism, which does not even in its
own thought find a firm ground, but constantly keeps
to the form of the question, is nevertheless not Montaigne's last
word. Pascal already saw this, when in the
masterly characterization of Montaigne found in his
``Pensées'' he sets it as the latter's aim to find
a morality founded on reason, without the light of faith.
Montaigne is a moralist, not merely because he psychologically
portrays the nature of man, but also because he
really expresses a coherent ethical view.
But of what kind can this be, when he takes up a
skeptical position toward every positive principle?

We have already reached the point where skepticism
in Montaigne passes over into a positive view. He
asserts ignorance, to be sure, and makes Que sais-je? his
principle; but ignorance is for him only a means of
grasping Nature. Men have, he says, by their
reasonings distorted the true and original nature,
whose traces are found only in the uncultivated and in the animals.
Nature has thereby lost its firm and universal
character and has received a particular form in each individual.
Virtue consists not in the dominion of thought or of reason,
but in letting nature develop according to its own
law, so that reason merely cuts off the false excrescences.
The drives follow their own natural course, and
reason cannot directly hinder them in this. Montaigne
% --- p. 218 ---
\setcounter{footnote}{0}%
\opage{218}therefore says that if he possesses any virtue, it
is an innocence that he owes to fortune rather than
to reason. When Nature of itself brings the ethical
harmony with it, it follows of itself that enjoyment
does not stand in opposition to virtue, but precisely through virtue
attains its highest and purest form. By nature we strive
for nothing but enjoyment. But the lower enjoyment
does not deserve this name; it is merely momentary
and full of anxiety and want. The higher enjoyment consists
in virtue itself, which one must therefore not, with the Stoics,
imagine as dark and austere, enthroned upon a steep
cliff in a lonely desert full of dangers. Nor
is it the case that virtue itself is mild and
smiling, but the way to it dark and arduous; then
no one would possess it, since no one manages to traverse
the whole way. The way must correspond to the goal to which it
leads. Virtue dwells on a fair and flowering height,
from which it sees all things beneath it, and one reaches it
by shady and flower-decked paths, with Nature
as guide, happiness and enjoyment as companions.

This doctrine recalls Epicureanism. But Montaigne
himself has already remarked that it is only a false
opinion that Epicureanism should fall short in firmness and severity
of Stoicism (liv. II chap. 11). If one
applies this to his own doctrine, it may be expressed
thus: that much labor and much
exertion are required for true nature to unfold in
man. What makes the Stoics stand further from Montaigne
is the ascetic element in them and their
disdain for the immediately natural and human,
which Montaigne precisely seeks to coax forth and to protect.
The greatness of the mind, he says, consists not so much in
striving upward and forward as in ordering and limiting itself.
Nothing is so fair and right as to be a man
in a good and proper way, no science so difficult as
to understand how to lead one's life in the right way, and the worst
% --- p. 219 ---
\setcounter{footnote}{0}%
\opage{219}sickness is contempt for our own being\footnote[1]{Entre nous --- siger han i Anledning af de kristelige Asketer --- ce sont choses que j'ay tousjours veues de singulier accord, les opinions supercelestes, et les moeurs soubterraines. Liv. III chap. 13.}. Nature is
a good guide, and no more good than wise and just.
Hence Montaigne nevertheless admires ``the loftiness of heroic
souls.'' It is indeed his principle that each shall be measured
by his own ell, according to his own nature (sur son propre
modele). His own weakness does not make him blind to
the strength of others; although he himself creeps on the ground, he can
still see the high flight of others, and he does not regard it
as any small matter to preserve the moral power of judgment
sound in an age like his, in which virtue does not
seem to be anything but ``un jargon de college,'' and
in which most strive to disparage every great and noble
action by subtly tracing it back to
base motives, so as to be able to draw it down into their own
sphere. Those ``heroic souls'' are models which
Nature selects in order to show to what height human
virtue and power can attain (liv. I chap. 36).

In the morality to which Nature leads us, enjoyment
thus attains its noblest form by becoming one with
virtue. It reconciles at the same time another opposition,
between egoism and devotion to others. By engaging
in public affairs and concerning oneself with
the affairs of others, one easily loses something of one's own being. Most
give themselves up immediately, and therefore do not belong to
themselves. Even sensual pleasure one must not give oneself wholly
up to; it hides pain at the bottom. One must not
immerse oneself in the world, but only touch it, glide
over its surface. The greatest thing in the world is to understand
how to belong to oneself. One must not bind oneself to anything
but one's own nature (n'espouser que soi-mesme). From this
it does not follow that one is to be an egoist and not care
% --- p. 220 ---
\setcounter{footnote}{0}%
\opage{220}about others; on the contrary, he who does not live for others
does not live for himself either. Regard for oneself includes
regard for others; both follow from the same nature.

Montaigne does not wish to set forth a slack and effeminate
doctrine of enjoyment. This is seen not only from what has gone before,
but also from the role the thought of death plays in
him. No one, he says, could rejoice in life
if he feared death, any more than the condemned
criminal on the way to the place of execution can rejoice in
the fair regions he passes through. Most people therefore
seek the expedient of trying to forget it, but this expedient
is bestial, is brutal stupidity. And besides, what
good does it do? One goes and comes, dances and laughs, as
if death did not exist; and when it then robs us of one
of our nearest, this carelessness is changed into the
deepest despair and dejection. Truly, even
if this bestial thoughtlessness were possible in a sensible
man, it would still be bought too dear. One must
in good time learn to fight this inevitable enemy,
to face it undaunted. Its most terrible weapon
is the strange and suddenly threatening. Let us then rob
it of this by making ourselves familiar with it, accustoming
ourselves to it by always having it in mind, by presenting
it to our imagination in all possible
shapes. Let us do as the ancient Egyptians, who
at their merriest feasts had mummies of dead
men brought in to serve as a reminder to the guests.
Since it is uncertain where death awaits us, let
us await it everywhere. To think about death is to think
about freedom; he who has learned to die has learned to cease
to be a slave; there is nothing evil in life for him who
has understood that the deprivation of life is not an evil. One
must not attach oneself so firmly to life, nor to any
enjoyment or deed in life, that one cannot part from it
without sorrow. Life is in itself neither a good
nor an evil; it depends on how one has
% --- p. 221 ---
\setcounter{footnote}{0}%
\opage{221}used it; one may have lived much in a short time. ``Je
veux, qu'on agisse et qu'on alonge les offices de la vie,
tant qu'on peult, et que la mort me trouve plantant
mes choulx, mais nonchalant d'elle et encore plus de
mon jardin imparfaict''. Besides, how our life is only
a vanishing nothing when we compare it with eternity,
or even merely with the age of the stars, the mountains,
the rivers and the trees! What right have we to
ascribe to ourselves a being in this brief moment, which is only
a little flash in the infinite course of an eternal night, so
short an interruption of our lasting and natural state,
since death takes up everything that lies before and after
this moment, indeed even a good part of it itself?
With right, therefore, Nature says to us: ``Go out of this
world as you came into it. The passage
you made from death to life, without passion and without
terror, you shall now make again from life to death. Your death
is a link in the order of the universe, in the life of the world''\footnote[1]{Liv. I chap. 19: Que philosopher c'est apprendre à mourir.}.

The character that man develops in the course of his
life springs of necessity from his original
nature. Conscious reason can, as we have already
heard, according to Montaigne act only indirectly, by warding off
and purifying; the positive fullness of life arises from the bosom
of Nature. Yet there is one positive feeling that reason can bring about,
and which arises when the effect of the infinitely
changing natural drive comes into conflict with the rational
insight into what is right. This inner disharmony
is felt as repentance. Thereby an inner tribunal is set up
in us, which no one can flee, even if he
appears toward others ever so calm and sure with regard to
his actions. Even if our deeds look brilliant
to others, they are for ourselves only like
clear rays that leap up from an impure and muddy
ground. But although repentance is a positive feeling, it
% --- p. 222 ---
\setcounter{footnote}{0}%
\opage{222}has its limit where reason has its; the very
passion that is one with our own being, and which
by habit and practice has taken still deeper root, it
cannot change. In this sense repentance is impossible.
It is only illusion or hypocrisy when one believes that one can truly
become, or have become, another. It
may happen at most, as Pythagoras thought, that
men receive a new soul when they approach the gods
to receive their oracles; but this new soul is
only assumed and vanishes again after a short time.
Everyone has in himself a form peculiar to him, governing his being
(une forme sienne, une forme maitresse),
which cannot be changed, but withstands every
external influence and transforms it in its own way. Therefore
the attempts to reform and improve the world reach only
the shells, never the kernel itself. As for myself,
says Montaigne, I can wish that I were
another man; I can condemn and abhor my general
nature (ma forme universelle) and pray God
for a complete reformation of my being; but this
I do not call repentance, any more than I call it
repentance to be dissatisfied at not being an angel or
a Cato. Because I imagine higher and more perfect
natures than my own, my faculties are not therefore improved,
any more than my arm or my understanding
becomes stronger because I can form an idea
of another in whom they are stronger. In that case
our most innocent actions too would become
an object of repentance, because we easily see that, on the assumption
of a perfect nature, they too would have
been carried out with greater perfection and dignity\footnote[1]{Liv. III chap. 2: Du repentir.}.

Thus: real repentance, repentance in the strictest
sense, would lead to my becoming another; but
% --- p. 223 ---
\setcounter{footnote}{0}%
\opage{223}another I cannot become so long as I am myself.
It is this contradiction in the concept of repentance
on which Montaigne bases his denial of its
reality. It is seen clearly here what role the principle of
immanence plays here in him, although he cannot formulate
it philosophically. In contrast to the orthodox
view of the corruption of human nature and
the necessity of an absolute rebirth, he asserts
the unity and immutability of human nature and
its connection with Nature as a whole. When
actions have been performed, they enter as necessary
links in the great process of the universe, in the infinite
series of causes and effects; and it is in vain to
grieve over them. Fancy cannot with all its strength
make a single point within this great process into nothing.
--- Such a view would, as e.g. Schopenhauer's
doctrine shows, be able to lead to pessimism; if one
namely assumes that the fundamental essence (la forme maistresse)
is evil. But Montaigne hesitates not a moment
over whether to give Heraclitus or Democritus the right: I
do not believe, he says, that there is so much misfortune among
us as there is vanity, or so much wickedness;
as there is folly!

Pascal has portrayed Montaigne as skeptic and
moralist, and on the whole as the one who presents ``the natural man.'' But he has not carried the portrayal
further in such a way that the ultimate foundation of Montaigne's
view of life came to light, the point on which everything
finally turns, and by which Montaigne reaches beyond
the merely popular and moral and indicates a whole philosophical
view of the world. This was, in and of itself, bound to have
less interest for Pascal, who himself took up an
avowedly practical standpoint. Pascal admires Montaigne's
dialectical carrying through of skepticism and accepts
it in the service of paradox; but he rejects his
morality as paganism, indeed in the last instance as cowardice.
% --- p. 224 ---
\setcounter{footnote}{0}%
\opage{224}Whether one can really charge with cowardice a view of life that makes it an essential task to practice the
thought of death and to have it as its confidant in everything one undertakes is another matter; in his burning passion
for the Gospel Pascal had to regard as too low everything that did not go up absolutely in this supernatural power.
That is also why he stops at the moral and religious consideration and does not advance to the metaphysical. It
likewise follows of itself that the metaphysical lies rather hidden in a causeur such as Montaigne. But after all
that we have heard in what precedes, it cannot be doubtful that behind his skepticism and his eudaemonism there rises
an infinite background, to which he constantly points, directly or indirectly. It is the idea of Nature in her
loftiness and infinity, which in its great chain also encloses human life as a single link. Sainte-Beuve, in his
Port-Royal (3e éd. II p. 412), where on the occasion of Pascal he comes to deal with Montaigne, has pointed out how
Montaigne, through his concept of nature, joins the fundamental thought of the great metaphysical systems, especially
Hobbes and Spinoza. The opposite of dualism, trust in human nature and its powers, had necessarily to lead a clear and
lively mind to the foundation on which the whole modern world view rests, even if he has not expressed this in a
philosophical formula.

The concept of nature already earlier made certain connections between Montaigne's skepticism and his morality
intelligible to us. But this concept plays a leading role with him in small matters as in great. If the talk is, e.g.,
of illness, he will have nothing to do with the physicians; they only botch Nature and hinder her in her own proper
development; --- like every other being, illness too has its determinate time to develop and to complete its life, and
it is in vain to intervene
% --- p. 225 ---
\setcounter{footnote}{0}%
\opage{225}here. One must let Nature rule; she understands things better than we. One must acquiesce in the order of
Nature. Like the harmony of the world, our life is composed of opposites, of different tones, which all belong in the
great coherence, sickness and death as well as health and life. --- If the talk is of education, it is likewise a
matter of letting Nature have her free course: nature peut tout et fait tout! --- Finally, by means of the concept of
nature he does away with his skepticism itself. The only way, he says, in which one can secure oneself against measuring
things with a false yardstick is to go beyond the relative comparisons and to set before one's eyes ``our mother Nature
in all her majesty''; she shows us a thoroughgoing and constant manifoldness, within which we discover ourselves and
everything that we call great; it is thereby that we first get the right concept of the magnitude of things. Without
such an absolute standard everyone regards his own religion and his own country's constitution as the only right ones
and calls everything else barbarism, whereas it ought rather to be called barbarism to narrow the power of ``our great
and mighty mother Nature'' by drawing arbitrary and accidental boundaries and barriers. Besides, it is also impossible;
Nature mocks our arts and triumphs over our ingenious inventions\footnote[1]{Ce n'est pas raison que l'art gaigne le poinct d'honneur sur nostre grande et puissante mere nature. Nous avons tant rechargé la beauté et la richesse par nos inventions, que nous l'avons du tout estouffee: si est ce que partout ou sa pureté reluict, elle faict une merveilleuse honte à nos vaines et frivoles entreprises. Liv. I chap. 30: Des cannibales.}. --- And with this concept of nature Montaigne's concept of God is,
finally, in close connection. ``Of all human opinions about religion'', he says (liv. II chap. 12), ``that seems to me
to have the most probability and justification which recognized God as an incomprehensible power, the source and sustainer
of all things, absolute goodness and perfection, who accepts and
% --- p. 226 ---
\setcounter{footnote}{0}%
\opage{226}approves the honor and the worship which men, each in his own way, rendered him, under whatever form, under
whatever name and in whatever manner it happened.'' We see from this how close a connection there is between
Montaigne's world view and his humane and tolerant intellectual bent.  Toleration is grounded with him not merely on
skepticism, but the presupposition that the eternal and infinite Nature underlies and reveals itself in all things and
all men must necessarily lead to the recognition that no one is cut off from the eternal and the divine, and that this
can come to light in countless forms, each of which has a right to develop with freedom.

Montaigne was not, as Pascal in his Jansenist zeal believed, an apostle of cowardice and sensuality, but a noble and
healthy nature. He mirrors himself in his writing, both in his good and his bad sides, with amiable naïveté. What
inwardness dwelt in him is seen, among other things, from the way in which he depicts his friendship with Etienne de la
Boétie\footnote[1]{Nos ames ont charié si uniement ensemble; elles se sont considerees d'une si ardente affection, et de pareille affection, descouvertes jusques au fin fond des entrailles l'un de l'aultre, que non seulement je cognoissois la sienne comme la mienne, mais je me feusse certainement plus volontiers fié à luy de moy, qu'à moy.} (liv. I chap.
27). This relation consisted, he says, in a complete mutual absorption of the one's will in the other's, so that no one
kept anything for himself, and each knew the other's innermost motives\footnote[2]{Etienne de la Boëtie (1530---1563), member of the Parliament of Bordeaux, is known especially for his ``Discours sur la servitude volontaire'', also called ``Contre un'', which he wrote in his eighteenth year (par maniere d'exercitation seulement, says Montaigne), and in which with glowing enthusiasm, though in a rhetorical style strongly reminiscent of the ancients, he proclaims the doctrine of the natural freedom of all and shows that it is
\opage{227}humanity's own fault when it is oppressed, for the oppressor is only one against the many and would have no power if these did not themselves submit to him. Montaigne calls him the greatest man of the century; later it has been thought that the beautiful chapter in which Montaigne speaks of his departed friend has made the latter more famous than his own writings would have done.}. ``When I compare
% --- p. 227 ---
\setcounter{footnote}{0}%
\opage{227}the whole rest of my life with the four years in which it was granted me to enjoy his dear company and
intercourse, it stands before me only as a fog, as a long and tedious night''. Friendship was on the whole the ethical
relation that suited Montaigne most; here he gave himself up entirely without fearing to lose himself. --- How healthy a
nature he was one sees from the remark that prosperity always puts him in a serious and thoughtful mood. ``I am more
devout in prosperity than in adversity. Fortune serves me for discipline and education, as misfortunes and chastisements
do for others; good luck is to me a special spur to moderation and modesty''. (Liv. III chap. 9).

Of necessity we had to return to the person himself, after apparently having left him in order to present his view of
life. In Montaigne less than in anyone else could the two be portrayed each by itself. But if we take the whole
phenomenon, it is a characteristic type of the Renaissance that we have before us here. Full-blooded and sensual, and
yet in his way directed toward the infinite, moving everything, and yet resting and enjoying in the midst of the
whirl of movement, delighting in antiquity and yet with the thought of the future in his lap. In his own time Montaigne
was understood only by a few\footnote[1]{See \emph{Nisard}: Précis de l'histoire de la littérature française. p. 211 f.}.
Justus Lipsius called him the French Thales, Cardinal du Perron called his book le bréviaire des honnétes
gens. In the seventeenth century Montaigne stood for Pascal
% --- p. 228 ---
\setcounter{footnote}{0}%
\opage{228}and the Jansenists as the great heathen who more clearly than anyone else has expressed the natural, i.e. the
sinful man's, principles. Only in the eighteenth century could he find kindred minds. Voltaire descends from him by his
fine and supple irony and sharp criticism. Rousseau expresses what Montaigne only hints at, the necessity of leaving the
refined life of culture (Montaigne's ``nature sophistiquée et raffinée'') and returning to genuine, original Nature.
Diderot, finally, unites both elements in himself and is therefore perhaps the later author who bears the greatest
resemblance to Montaigne.

We have tried to find a system in Montaigne's Essays; and with the system there must necessarily follow a criticism.
As we have sought to show, Montaigne's main thought is the same that presses forward in the whole thinking of the
Renaissance, the idea of Nature. This idea appears with him in a peculiar form. It forms the general background, the
last presupposition of his skepticism; it is likewise the positive power that he sees at work in everything, and for
which he everywhere seeks to prepare a place. But there is nevertheless, especially when we consider him as an
ethicist, a decisive defect in his conception of the concept of nature. He sees the natural as an immediately
dissolving stream, as a kind of vegetating, which is hindered and distorted by conscious thinking and acting.
From this follows in turn that man's thinking and acting receive only a negative or indirect significance, instead of
these functions, just as much as any other, having to have their place and significance within Nature. The principle of
human action must be found precisely by investigating what application of these functions follows from man's true
nature. They thus do not stand outside the stream of infinite Nature, but are themselves forms of it. If this
consideration is carried through, the active element in the concept of nature steps forth more clearly; in place of the
immediate flowing there now steps
% --- p. 229 ---
\setcounter{footnote}{0}%
\opage{229}the conscious devotion and the energetic development of the task indicated in the idea of Nature. From this
side we find the matter in part conceived in Giordano Bruno, Montaigne's younger contemporary, when, like Plato, he sees
the natural drive as the expression of an ideal unrest, which does not allow man to rest until he has acquired harmony
and reconciliation through the labor of thought and will and through devotion to the infinite with his whole
personality\footnote[1]{From this there follows, of course, also a different conception of the decisive ethical problems. Thus Bruno in particular maintains the reality of repentance, which Montaigne in a naturalistic fashion denied. But Bruno conceives repentance in such a way that he does not give up the principle of immanence. He assumes no absolute evil, no corruption of human nature. But without strife and struggle, without contradiction and resistance the true and the good cannot prevail: and evil is the negative power that raises the contradiction and makes the struggle necessary. This negative power must be sublated; the soul wins itself through contrition. Repentance is thus a sign of the ideal nature of the soul: it is like a shining spark that strikes out of the hard flint and bears witness to its kinship with the sun. --- It is clear that what is of principle and justified in Montaigne's polemic is not denied by Bruno's conception; for that polemic was directed against repentance as a complete transformation of nature.}.
In Bruno there is therefore found no such opposition as in Montaigne between finite thinking and will, which can accomplish
nothing, and the view of infinite Nature. Montaigne's skepticism, which seen from another side is one of his chief
merits, in that it frees us from the presuppositions of faith in authority, has nevertheless precisely its ground in
the fact that the idea of Nature with him, however much it is the basis of everything, stands as a distant and mystical
power and is not drawn forth to concrete realization in reality.

But it is of course a one-sided way of looking at things to see Montaigne as a philosophical ethicist. He comes into
his own only when one sees him in his whole
% --- p. 230 ---
\setcounter{footnote}{0}%
\opage{230}literary-historical significance, and this rests on the fact that he does not address himself to the learned and the
solitary thinkers, but that in a style so light, clear and lively as had not yet been heard in a modern language, he
directed his speech to everyone who would hear him, if for no other reason, then at least to amuse himself in so
agreeable a company as Montaigne's Essays are.
\signature{\textbf{Harald Høffding.}}

\end{document}
